% Four replacements in Appendix A.7. Lines 2264-2335 (payoffs over the reserve domain,
% the continuation object, (A.42)-(A.44)) stand without change. No \NUM placeholders:
% every value below is a formula, and the endpoint revenue is \rho h/2.
% Removed here: Proposition A.10 and equations (A.45)-(A.49). The integrator decides
% whether later equation tags keep (A.50) onward or close the gap. These fragments
% keep (A.50) and (A.51) so that no other reference moves.

% REPLACES: main.tex lines 2336-2465 ("\subsubsection{Price pooling at fixed orders} ..." to "... and revenue fall as the cutoff rises.")
\subsubsection{The price experiment at fixed
orders}\label{the-price-experiment-at-fixed-orders}

On the reserve domain \(p\in[0,h]\) gross acquisition profit is
nonnegative for each challenger value, so it is nonnegative at every belief.
The zero-cost realization therefore prepares at every price, strictly
whenever gross profit is positive and by the tie rule when it is zero.
Entry is at least \(\rho\) at every price and at every reserve in this domain.
Preparation cannot vanish on a set of flows, so no set of flows is left
to collect at a no-entry price.

Fix the orders and a reserve with \(\Delta_T(p,r)=t_H(p,r)-t_L(p,r)>0\).
Write \(B_{p,r}(\mu)\) for belief-weighted gross profit at that reserve and
\(e(\mu)\) for the entry probability it induces under the two-point cost law.
Then \(e(\mu)\ge\rho\), \(e\) is nondecreasing, and competitive pricing gives
\(P(\mu)=t_0(p,r)+e(\mu)[t_L(p,r)-t_0(p,r)+\Delta_T(p,r)\mu]\).
The bracket is nonnegative. A prepared low-value challenger that meets the reserve
pays \(\max\{p,\min(R,\ell)\}\ge p\mathbf1\{R\ge p\}\) in each incumbent state, and
one that fails the reserve leaves the no-challenger payment in place, so
\(t_L(p,r)\ge t_0(p,r)\).
For \(\mu_2>\mu_1\) the price increment is therefore at least
\(\rho\Delta_T(p,r)(\mu_2-\mu_1)>0\).
Distinct raw order-flow posteriors give distinct prices, and the buyer
recovers \(\mu_X\) from the price. This is the increment argument of
\hyperref[result-proposition-a-2]{Proposition A.2} read at a general reserve.
Only flows that share one posterior share a price. A Laplace plateau is such a
set, and its price is then an atom with that common posterior.

Two continuations with the same orders and a positive spread therefore have the
same price mapping. They remain different elements of \(\mathcal E(p,r)\) only
through different orders. The seller's problem \eqref{eq:paper-a-44} still needs
a selection after every reserve, because the order profile itself can vary with
the reserve and can be nonunique at one reserve.
Dow et al. (\citeproc{ref-DowGoldsteinGuembel2017}{2017}) distinguish
continuations by their pooling regions in a setting without this entry floor.
That margin is absent here on the part of the domain with a positive spread.
At a reserve with \(\Delta_T(p,r)=0\) the price carries no spread to invert, and
the constant-price candidate is checked directly. No comparison below uses such
a reserve.

% REPLACES: main.tex lines 2466-2512 ("\subsubsection{Exact reserve events} ..." to "... comparison.")
\subsubsection{Exact reserve events}\label{exact-reserve-events}

One participation equality remains in the interior of the reserve domain, the
expensive-entry event. The zero-cost equality sits at the endpoint of the
domain. Neither event is an interior point of a strict-inequality region, and
the same continuation need not be supported on both sides of the interior event.

\emph{Domain endpoint at the high value.} When the reserve exceeds incumbent support
and excludes the low class, \(g_H=h-p\), \(g_L=0\), \(t_H=p\), and
\(t_0=t_L=0\). At the lowest feasible belief gross profit is \(B_p(m)=m(h-p)\).
It is strictly positive for every \(p<h\) and equals the zero preparation cost only at

\[
p_L=h.
\tag{A.50}\label{eq:paper-a-50}
\]

At that endpoint the high-value challenger earns no rent. The tie rule admits
the zero-cost realization at every price, with equality \(B(\mu)=0\) at each
belief; the high-cost realization stays out because \(0<c_H\).
Entry is therefore \(\rho\) at every price. The inversion above and the global
bound of \hyperref[pa-proofs]{Appendix A.2} apply with \(e\ge\rho\) under
equality, giving unique full orders. Then \(\mathsf E=\rho\), the sale
probability is \(\rho/2\), and \(\mathcal R_T=\rho h/2\). Preparation
here never produces two admissible bidders: the incumbent is excluded.
Offsets from this event run only below \(h\), inside the domain, where the
zero-cost comparison is strict. There is no interior low-cost floor event left
to locate.

\emph{Strong-incumbent expensive entry.} With the low class excluded and
\(\ell<p<r\), \(g_H=h-r/2-p^2/(2r)\) and \(g_L=0\). The upper-posterior
entry event \(Mg_H=c_H\) occurs at

\[
p_H=\sqrt{2r\left(h-\frac{c_H}{M}\right)-r^2},
\tag{A.51}\label{eq:paper-a-51}
\]

which at \(r=r_1\) lies in \((\ell,r_1)\). There \(\tau=M\) and
\(x^*=1\), so \(\alpha_H=\tfrac12\) and \(\alpha_L=\tfrac12e^{-2/b}\):
the plateau has positive probability and entry at indifference admits
the expensive type on it. The zero-cost realization prepares strictly there,
since at the lowest belief \(mg_H=mc_H/M>0\). No separate floor condition
is needed. The global full-order bound is checked
separately from this equality. Both events are evaluated from their
defining identities, not from a rounded reserve and a floating-point
comparison.

% REPLACES: main.tex lines 2527-2536 ("the preparation, admissible-challenger, sale, ..." to "... rather than as e_H/2.")
the preparation, admissible-challenger, sale, and two-admissible-bidder
probabilities. Inclusion and exclusion give
\(\mathsf S=\Pr(R\ge p)+\mathsf A-\mathsf C_2\), and
\(0\le\mathsf C_2\le\mathsf A\le\mathsf E\le1\). Preparation is not the
same event as two admissible bidders. At the endpoint event \(p_L=h\),
\(\mathsf C_2=0\) while \(\mathsf E=\rho>0\). At the strong-incumbent ceiling
event, \(p_H<r_1\) and \(\mathsf C_2=(1-p_H/r_1)e_H/2>0\), because a
prepared high-value challenger and the incumbent can both be admissible.
Outside the benchmark support, high-value ownership is computed from the
actual allocation event rather than as \(e_H/2\).

% REPLACES: main.tex lines 2564-2583 ("The term -p/r is the direct effect ..." to "... remains open.")
The term \(-p/r\) is the direct effect on acquisition profit; the
posterior derivative is the change in information generated by trading,
which vanishes locally when orders stay at their bounds and can be
nonzero when they adjust. The decomposition requires a continuously
differentiable pure branch, a continuously differentiable cost CDF, and
an integrable dominating function for the integrand in
\eqref{eq:paper-a-54}; Online Appendix A.10 states these conditions.
There is no general sign restriction on the information term. At an
atomic preparation threshold or on a nonregular
continuation, the level objective \eqref{eq:paper-a-44} with the
complete conditional laws is the correct object.

The reserve sweep in Online Appendix C.6 records continuations found
under the declared searches, with continuation identity taken from the
full object \eqref{eq:paper-a-43} and with the endpoint reserve
\eqref{eq:paper-a-50} and the ceiling event \eqref{eq:paper-a-51} included. It does not
prove existence at unresolved reserves, exhaust \(\mathcal E(p,r)\), or
establish a global seller optimum. Characterizing seller-optimal terms,
their commitment timing, and whether they preserve the entry reversal
remains open.
